%======================================================================
\chapter{Discussion, Limitations and Conclusion}
\label{ch:conc}
%======================================================================

\section{Where our implementation differs from the base paper}
\label{sec:deviations}

Reporting an implementation as though it reproduced a paper exactly, when it did not, would
be misleading. Four differences are material and we set them out explicitly.

\begin{table}[H]
\centering
\caption{Deliberate divergences from the base paper.}
\label{tab:deviations}
\small
\begin{tabularx}{\textwidth}{@{}L{6mm} L{34mm} L{40mm} X@{}}
\toprule
& \textbf{Paper specifies} & \textbf{We implemented} & \textbf{Reason and effect}\\
\midrule
V1 & Pentatope ECC (PEECC) over a 5D point group, claimed DLP complexity $\approx 2^{5n/2}$
   & Standard secp256r1 (NIST P-256), 2D, $\approx 2^{2n/2}$
   & The paper gives no concrete curve parameters, no reference implementation and no test
     vectors for PEECC, so we could not reproduce it with confidence. Our results are
     therefore a \emph{lower bound} on the paper's claimed security level, not a
     reproduction of it. Everything else in the protocol is independent of this choice.\\
\addlinespace[3pt]
V2 & $E\{\PU_{recv},\langle m \rangle\}$ --- direct public-key encryption; AES-128 named in
     the energy analysis
   & ECDH key agreement followed by AES-256-GCM
   & This is the standard hybrid realisation of public-key encryption (the construction
     ECIES uses). It is functionally equivalent for the protocol's purposes and is what a
     real deployment would do, since direct public-key encryption of every message would be
     prohibitively expensive.\\
\addlinespace[3pt]
V3 & Adaptive $\Delta T$ from a rolling average of observed inter-node delays
   & Fixed $\Delta T = 5$\,s
   & The adaptive mechanism is described but not specified numerically. A fixed window is
     the conservative choice and is what our replay test exercises. Making it adaptive is
     listed in the future work.\\
\addlinespace[3pt]
V4 & Signature-based validation at registration (Proposition~1)
   & Hash-based verification value $\vartheta_i = H(\RID_i \cat \PU_i)$
   & We implemented the registration formulae of Section~IV-B, which are hash-based. The
     paper's Proposition~1 additionally invokes a signature; the two descriptions in the
     paper are not fully aligned, and we followed the concrete formulae rather than the
     prose.\\
\bottomrule
\end{tabularx}
\end{table}

\section{Limitations}

\begin{itemize}
\item \textbf{Simulation, not hardware.} Energy figures come from published ARM Cortex-M4
      benchmarks, not from measurements on an acoustic modem. Real modems (EvoLogics S2C,
      Teledyne Benthos ATM-900) have transmit power in the watts, so the TX and RX figures
      here are optimistic by orders of magnitude in absolute terms --- though their
      \emph{ratio} to the cryptographic cost, which is the point being made, holds.
\item \textbf{Independent loss.} Our Bernoulli model treats losses as independent. Real
      acoustic channels have bursty loss driven by ship noise, thermoclines and surface
      conditions, for which a Gilbert--Elliott two-state model would be more faithful.
\item \textbf{Simplified mobility.} The $\pm10$\,m per-round drift is one-dimensional, and
      does not model depth changes, currents, or nodes moving out of range entirely.
\item \textbf{Scyther is bounded.} Verification is at \code{max-runs=5}. This is standard
      practice and no attack was found, but it is a bounded search, not an unbounded proof.
\item \textbf{The anomaly detector is offline.} \code{detect\_anomalies} runs over the
      accumulated delay log at the end of a run. A deployed version would need a streaming
      estimator with an exponentially weighted moving average.
\item \textbf{No revocation implemented.} The paper's revocation list is described in
      Chapter~\ref{ch:proto} but our simulator does not implement or exercise it.
\end{itemize}

\section{Future work}

\begin{enumerate}
\item \textbf{Post-quantum key agreement.} Shor's algorithm breaks ECDH on a sufficiently
      large quantum computer, and underwater deployments have service lives measured in
      years. Substituting CRYSTALS-Kyber (NIST's selected KEM) for ECDH is tractable
      precisely because of the protocol's hop-wise modularity: only the key-agreement layer
      changes, and the message structure, freshness logic and fallback all stay as they are.
      The cost is message size --- a Kyber-768 ciphertext is 1088 bytes against 32 for an
      ECDH point --- which on a 10\,kbps channel is the whole difficulty and would need
      study.
\item \textbf{Implement PEECC.} Building the pentatope construction with concrete parameters
      would let the paper's $2^{5n/2}$ complexity claim be tested rather than assumed, and
      would close divergence V1.
\item \textbf{Adaptive $\Delta T$.} Maintain a rolling mean $\mu_d$ and standard deviation
      $\sigma_d$ of observed inter-node delay and set $\Delta T = \mu_d + 3\sigma_d$. This
      would tighten the replay window in calm conditions and widen it under adverse ones,
      improving both security and reliability over a fixed 5\,s.
\item \textbf{Hardware benchmarking.} Measure real energy on an acoustic modem, replacing
      the literature figures with instrumented ones.
\item \textbf{Hybrid acoustic--optical links.} Underwater optical modems reach Mbps rates
      below about 100\,m. A protocol that selects the link type by distance would transform
      the throughput picture for short hops such as UWS $\to$ SUB, where 150\,m is well
      within optical range.
\item \textbf{Multipath routing.} Rather than failing over to a backup only after the
      primary dies, send over multiple paths concurrently for redundancy --- valuable where
      a lost authentication means a lost sensing window.
\item \textbf{Streaming anomaly detection with response.} Move the $z$-score detector
      online and connect it to the Alerts field the base paper already carries in $M_6$,
      $M_9$ and $M_{12}$, so that a detected injection propagates to the base station.
\end{enumerate}

\section{Conclusion}

The base paper addresses a genuine gap: authentication for underwater acoustic networks,
where the channel is slow, narrow, lossy and energy-scarce enough that conventional security
is unusable. Its design decision --- authenticate each hop independently with a fresh nonce
and a bounded timestamp, rather than end to end --- is the right one for the medium, and it
is what makes the reported 2112-bit, 0.4\,ms, 48.8\,\uJ{} figures achievable.

We reproduced that protocol as a working eight-node simulator and verified it formally. In
doing so we corrected a serious flaw in our own first implementation: a repeating-key XOR
cipher, which offered no real confidentiality and no integrity at all, was replaced with
AES-256-GCM authenticated encryption. We replaced invented delay constants with the Thorp
absorption model, added a 15\,\% loss channel with retransmission, tracked battery
depletion per node, gave the AUVs motion, and built a detector for the delay-injection
attack the paper names but does not counter. We restructured the Scyther model from one
tangled five-role protocol into four independent Needham--Schroeder--Lowe hops, which is
both more faithful to the paper's design and what allowed all 32 claims to verify.

Two findings are worth carrying forward. First, in this setting \textbf{message size
dominates cryptographic cost}: transmit and receive account for 64\,\% of the per-hop energy
budget, so the paper's communication-overhead reduction is a more valuable contribution than
its computational one. Second, \textbf{formal verification found a real design problem} ---
our monolithic model's failure to establish \code{Niagree} was not a tooling artefact but a
symptom of nonces crossing hop boundaries, and fixing the model meant fixing the design.

The protocol as implemented authenticates a hop in 0.4\,ms and 48.8\,\uJ, sustains 98.65\,\%
end-to-end reliability over a channel that drops one packet in seven, reroutes around a
failed buoy without restarting the session, rejects replayed and tampered messages, and
carries a machine-checked proof that its message design admits no attack by a
Dolev--Yao adversary.
